% BHU PYQ -- Transcribed & Corrected
% Paper: CHEMJ21/CHEMN21: Basic Concepts of Chemistry-II
% Exam: B.Sc. (Hons) Semester II (NEP) Examination, 2024-25
% Ref Code: CECONF/N/2024-25/10 & 11
% Category: Major / Minor
% Department: Chemistry

\documentclass[11pt,a4paper]{article}
\usepackage[a4paper,top=2.5cm,bottom=2.5cm,left=2.8cm,right=2.8cm,headheight=14pt]{geometry}
\usepackage{amsmath,amssymb}
\usepackage[T1]{fontenc}
\usepackage[utf8]{inputenc}
\usepackage{lmodern,microtype,fancyhdr,enumitem,hyperref,lastpage,parskip}

\hypersetup{
    colorlinks=true,
    urlcolor=black,
    linkcolor=black,
    pdftitle={CHEMJ21/CHEMN21: Basic Concepts of Chemistry-II -- BHU Sem II 2024-25 (NEP)}
}

\pagestyle{fancy}
\fancyhf{}
\renewcommand{\headrulewidth}{0.4pt}
\renewcommand{\footrulewidth}{0.4pt}
\fancyhead[L]{\small \textit{Banaras Hindu University}}
\fancyhead[R]{\small \textit{CHEMJ21/CHEMN21: Basic Chemistry-II (2024-25)}}
\fancyfoot[C]{\small Page \thepage\ of \pageref{LastPage}}

\newcommand{\pts}[1]{\hfill \textit{[#1]}}

\begin{document}

\begin{center}
  {\large\bfseries BANARAS HINDU UNIVERSITY}\\[2pt]
  {\normalsize B.Sc. (Hons) Semester II Examination, 2024-25 (NEP)}\\[6pt]
  \rule{0.8\linewidth}{0.5pt}\\[6pt]
  {\large\bfseries CHEMISTRY}\\[4pt]
  {\large\bfseries Paper Code: CHEMJ21 / CHEMN21 : Basic Concepts of Chemistry-II}\\[2pt]
  {\small Ref. No.: CECONF/N/2024-25/10 \& 11}\\[4pt]
  {\normalsize Time: 3 Hours\hfill Maximum Marks: 70}
\end{center}

\rule{\linewidth}{0.4pt}

\smallskip

\noindent \textit{Instructions: Attempt five questions including Question 1 which is compulsory. Attempt at least one question from Questions 2 and 3. The figures in the right-hand margin indicate marks.}

\bigskip

\noindent \textbf{Question 1.} Answer all parts:
\begin{enumerate}[label=(\alph*)]
  \item ``Carbon possesses only two unpaired spins in its ground state yet it forms four bonds.'' Why? \pts{2}
  \item Comment upon stability of $\text{Be}_2$ as per molecular orbital theory. \pts{2}
  \item ``The HCF bond angle is slightly smaller than HCH in $\text{CH}_3\text{F}$.'' Explain this observation in terms of Bent's rule. \pts{1}
  \item Explain why the phenols are stronger acids than alcohols but weaker than carboxylic acids. \pts{3}
  \item Why fully-eclipsed Newman projection of n-butane have maximum Torsional Strain? \pts{2}
  \item Why does only the E2 elimination show the isotope effect and E1 does not? \pts{1}
  \item In Williamson's ether synthesis, why primary alkyl halide used instead of a secondary or tertiary one? \pts{2}
  \item How many transition states exist in the $\text{S}_\text{N}1$ reaction? \pts{1}
\end{enumerate}

\bigskip

\noindent \textbf{Question 2.}
\begin{enumerate}[label=(\alph*)]
  \item Explain in brief the Valence bond theory of metallic bonding through a suitable example. \pts{2}
  \item ``The boiling point of HF is much higher than that of HCl, HBr and HI while in aqueous state it is poorest acid among all.'' Why? \pts{4}
  \item Discuss the structure and bonding of either Methyl Lithium or diborane. \pts{4}
  \item Draw the crystal structure of Rock salt and mention the co-ordination nos. of sodium and chloride ions. \pts{2}
  \item Why is Si-F bond stronger than corresponding bond with Carbon? \pts{2}
\end{enumerate}

\bigskip

\noindent \textbf{Question 3.}
\begin{enumerate}[label=(\alph*)]
  \item Let A and B be two atoms with wave functions for 1s orbitals as $\psi_A$ and $\psi_B$ respectively. Write the molecular orbital wave functions using these atomic orbital wave functions $\psi_A$ and $\psi_B$. Write expressions using molecular orbitals for the probability of finding the electron density in a unit volume and interpret the meaning of each term in the expressions. \pts{5}
  \item ``The compounds trimethylamine $(\text{CH}_3)_3\text{N}$ and Trimethylsilylamine $(\text{Me}_3\text{Si})_3\text{N}$ have similar formulae but have pyramidal and triangular planar structures respectively.'' Why? \pts{3}
  \item Explain in brief the molecular orbital theory of metallic bonding through a suitable example. \pts{2}
  \item Comment upon variation of bond angles in $\text{CH}_4$, $\text{NH}_3$ and $\text{H}_2\text{O}$ in the light of VSEPR theory. \pts{2}
  \item Explain the defects noted in extending homonuclear diatomic molecular orbital concepts in the heteronuclear diatomic molecule CO using energy level diagram. \pts{2}
\end{enumerate}

\bigskip

\noindent \textbf{Question 4.}
\begin{enumerate}[label=(\alph*)]
  \item Write down the product(s) of the following reactions: \pts{8}
  \begin{enumerate}[label=(\roman*)]
    \item Cyclopentadiene + Maleic anhydride $\xrightarrow{0^\circ\text{C}}$ ?
    \item $\text{H}_3\text{C}-\text{C}\equiv\text{CH} \xrightarrow{\text{dil. }\text{H}_2\text{SO}_4 / \text{HgSO}_4}$ ?
    \item Cyclopentene $\xrightarrow{1.\text{ MCPBA} / 2.\text{ H}^+}$ ?
    \item $n\text{H}_3\text{C}-\text{CH}=\text{CH}_2 \xrightarrow{\text{Benzoyl peroxide} / \Delta}$ ?
  \end{enumerate}
  \item Arrange the following in increasing order of basicity with a suitable explanation: \pts{4}
  \begin{enumerate}[label=(\roman*)]
    \item $\text{Ph-NH}_2$
    \item $\text{Ph}_2\text{NH}$
    \item $\text{Ph}_3\text{N}$
  \end{enumerate}
  \item Predict the E and Z nomenclature of the following: \pts{2}
  \begin{enumerate}[label=(\roman*)]
    \item 1-bromo-1-chloropropene
    \item 2-bromo-1-chloropropene
  \end{enumerate}
\end{enumerate}

\bigskip

\noindent \textbf{Question 5.}
\begin{enumerate}[label=(\alph*)]
  \item Write down the major products formed in the following reaction with a suitable mechanistic explanation: \pts{8}\\
  $\text{CH}_3-\text{CH}_2-\text{C}(\text{CH}_3)(\text{Br})-\text{CH}_3 \xrightarrow{\text{KOH/EtOH}}$ Substitution (?) / Elimination (?)
  \item What is the intermediate involved in the E1cB mechanism? \pts{2}
  \item Identify the products formed in the following reaction with a suitable mechanism: \pts{4}\\
  $\text{CH}_3-\text{O}-\text{C}(\text{CH}_3)_3 + \text{HI} \rightarrow$ ?
\end{enumerate}

\bigskip

\noindent \textbf{Question 6.}
\begin{enumerate}[label=(\alph*)]
  \item Predict the R and S nomenclature of glyceraldehyde and threonine derivatives. \pts{3}
  \item How do you prepare the following from phenylmagnesium iodide? \pts{6}
  \begin{enumerate}[label=(\roman*)]
    \item Benzenesulfonic acid ($\text{Ph-SO}_2\text{H}$)
    \item Phenol ($\text{Ph-OH}$)
    \item Benzyl alcohol ($\text{Ph-CH}_2\text{OH}$)
  \end{enumerate}
  \item How do you prepare the following from ethyl acetoacetate? \pts{5}
  \begin{enumerate}[label=(\roman*)]
    \item 1,4-diketone
    \item 1,6-diketone
  \end{enumerate}
\end{enumerate}

\bigskip

\noindent \textbf{Question 7.}
\begin{enumerate}[label=(\alph*)]
  \item Identify the products formed in the following reaction with a suitable mechanism: \pts{4}\\
  $\text{CH}_2=\text{CH}-\text{CH}=\text{CH}_2 \xrightarrow{\text{Br}_2}$ ?
  \item How do you prepare the following from diethyl malonate? \pts{4}
  \begin{enumerate}[label=(\roman*)]
    \item Glutaric acid
    \item Dimethyl acetic acid
  \end{enumerate}
  \item Discuss the $\text{S}_\text{N}\text{i}$ reaction with a suitable example. \pts{4}
  \item Distinguish between primary and secondary alcohols by the substitution method. \pts{2}
\end{enumerate}

\bigskip
\begin{center}
  \textbf{***}
\end{center}

\end{document}
